\documentclass{exam-zh}
\usepackage{edu-practice}

\examsetup{
  question/show-answer = false,
  fillin/show-answer = false,
  solution/show-solution = hide,
}

\begin{document}

\section*{26一模相似三角形证明题精选}

% ---------- 第 1 题 (宝山23) ----------

\needspace{8\baselineskip}

\begin{problem}[points=12]
  \noindent
  \begin{minipage}[t]{\dimexpr\linewidth-58mm-6mm\relax}
    \vspace{0pt}
    如图，在菱形 $ABCD$ 中，$E$ 是 $CD$ 上一点，联结 $BE$ 并延长，交 $AD$ 的延长线于点 $G$，交 $AC$ 于点 $F$。
    （1）求证：$BE\cdot BF=BG\cdot EF$；（2）若 $E$ 是 $CD$ 的中点，且 $AF\cdot CF=2BF^2$，求证：$BE\perp CD$。
  \end{minipage}\hfill
  \begin{diagramcoltikz}{58mm}{}
  \includegraphics[width=\linewidth]{\detokenize{assets/Q001/manual-prompt-1-5d41322a0c38e9c2.png}}
  \end{diagramcoltikz}
  \par\medskip
  \answerarea[72mm]
\end{problem}

% ---------- 第 2 题 (崇明23) ----------

\needspace{8\baselineskip}

\begin{problem}[points=12]
  \noindent
  \begin{minipage}[t]{\dimexpr\linewidth-52mm-6mm\relax}
    \vspace{0pt}
    如图，在 $\triangle ABC$ 中，点 $D、E$ 分别在边 $BC、AC$ 上，$AB=AD$，$AD$ 与 $BE$ 交于点 $F$，且 $BC^2=AC\cdot EC$。
    （1）求证：$\angle ADB=\angle BEC$；（2）如果 $AB=BC$，求证：$AE\cdot AD=AF\cdot BE$。
  \end{minipage}\hfill
  \begin{diagramcoltikz}{52mm}{}
  \includegraphics[width=\linewidth]{\detokenize{assets/Q002/prompt-figure.png}}
  \end{diagramcoltikz}
  \par\medskip
  \answerarea[72mm]
\end{problem}

% ---------- 第 3 题 (奉贤23) ----------

\needspace{8\baselineskip}

\begin{problem}[points=12]
  \noindent
  \begin{minipage}[t]{\dimexpr\linewidth-52mm-6mm\relax}
    \vspace{0pt}
    如图，在 $\triangle ABC$ 中，$D,E$ 分别在 $AB,BC$ 上，$CD,AE$ 交于 $F$，且 $CD=BD,AE=AC$。
    （1）证明 $\triangle ADF\sim\triangle CDA$；（2）联结 $DE$，由题给比例证明 $AD\cdot AE=DC\cdot DE$。
  \end{minipage}\hfill
  \begin{diagramcoltikz}{52mm}{}
  \includegraphics[width=\linewidth]{\detokenize{assets/Q003/manual-prompt-1-821ff9e637f000da.png}}
  \end{diagramcoltikz}
  \par\medskip
  \answerarea[72mm]
\end{problem}

% ---------- 第 4 题 (嘉定23) ----------

\needspace{8\baselineskip}

\begin{problem}[points=12]
  \noindent
  \begin{minipage}[t]{\dimexpr\linewidth-52mm-6mm\relax}
    \vspace{0pt}
    如图10，在 $\triangle ABC$ 中，$\angle ACB=90^\circ$，点 $D$ 是边 $AB$ 上一点，满足 $AC^2=AD\cdot AB$，点 $E$ 是 $AC$ 的中点，联结 $ED$ 并延长交 $CB$ 的延长线于点 $F$。
    求证：（1）$CD\perp AB$；（2）$BC\cdot CF=AC\cdot DF$。
  \end{minipage}\hfill
  \begin{diagramcoltikz}{52mm}{}
  \includegraphics[width=\linewidth]{\detokenize{assets/Q004/manual-prompt-1-df43e2ffd7de8673.png}}
  \end{diagramcoltikz}
  \par\medskip
  \answerarea[72mm]
\end{problem}

% ---------- 第 5 题 (闵行22) ----------

\needspace{8\baselineskip}

\begin{problem}[points=12]
  \noindent
  \begin{minipage}[t]{\dimexpr\linewidth-58mm-6mm\relax}
    \vspace{0pt}
    如图，线段 $AD,BC$ 交于点 $E$，$F$ 是 $ED$ 中点，联结 $AB,BD,CD$，延长 $BA,FC$ 交于 $G$。
    已知 $\angle BAD=90^\circ,\frac{AE}{CE}=\frac{BE}{DE}$。
    （1）证明 $\angle ABE=\angle FCD$；（2）若 $DA$ 平分 $\angle BDC$，证明 $\frac{BG}{BD}=\frac{BC}{2CF}$。
  \end{minipage}\hfill
  \begin{diagramcoltikz}{58mm}{}
  \includegraphics[width=\linewidth]{\detokenize{assets/Q005/manual-prompt-1-5786aa6de12c1101.png}}
  \end{diagramcoltikz}
  \par\medskip
  \answerarea[72mm]
\end{problem}

% ---------- 第 6 题 (普陀23) ----------

\needspace{8\baselineskip}

\begin{problem}[points=12]
  \noindent
  \begin{minipage}[t]{\dimexpr\linewidth-52mm-6mm\relax}
    \vspace{0pt}
    已知：如图，四边形 $ABCD$ 中，点 $E$ 在边 $BC$ 上，$DE$ 交 $AC$ 于点 $F$，$\angle AED=\angle B$，$AE^2=AF\cdot AC$。
    （1）求证：$\triangle ABE\sim\triangle ECF$；（2）如果 $\dfrac{AD}{AE}=\dfrac{AC}{EC}$，求证：$DE\cdot AC=AE\cdot BC$。
  \end{minipage}\hfill
  \begin{diagramcoltikz}{52mm}{}
  \includegraphics[width=\linewidth]{\detokenize{assets/Q006/manual-prompt-1-2eb6a34e993af7e3.png}}
  \end{diagramcoltikz}
  \par\medskip
  \answerarea[72mm]
\end{problem}

% ---------- 第 7 题 (青浦23) ----------

\needspace{8\baselineskip}

\begin{problem}[points=12]
  \noindent
  \begin{minipage}[t]{\dimexpr\linewidth-52mm-6mm\relax}
    \vspace{0pt}
    如图9，在平行四边形 $ABCD$ 中，点 $E$ 在边 $AD$ 上，$BE、CD$ 的延长线交于点 $F$，联结 $CE$，已知 $CE^2=DE\cdot BC$。
    （1）求证：$\triangle ABE\sim\triangle EFC$；（2）求证：$CE\cdot CF=EF\cdot AD$。
  \end{minipage}\hfill
  \begin{diagramcoltikz}{52mm}{}
  \includegraphics[width=\linewidth]{\detokenize{assets/Q007/prompt-01.png}}
  \end{diagramcoltikz}
  \par\medskip
  \answerarea[72mm]
\end{problem}

% ---------- 第 8 题 (松江23) ----------

\needspace{8\baselineskip}

\begin{problem}[points=12]
  \noindent
  \begin{minipage}[t]{\dimexpr\linewidth-52mm-6mm\relax}
    \vspace{0pt}
    如图，在梯形 $ABCD$ 中，$AD\parallel BC,AB\perp BC$，$E$ 在 $AB$ 上，$DE$ 与 $AC$ 交于 $G$，$CE$ 平分 $\angle ACB$ 且 $DE\perp CE$。
    （1）证明 $AE^2=AG\cdot AC$；（2）证明 $AG\cdot GC=EG\cdot ED$。
  \end{minipage}\hfill
  \begin{diagramcoltikz}{52mm}{}
  \includegraphics[width=\linewidth]{\detokenize{assets/Q008/manual-prompt-1-e755d287c9f78d2c.png}}
  \end{diagramcoltikz}
  \par\medskip
  \answerarea[72mm]
\end{problem}

% ---------- 第 9 题 (徐汇23) ----------

\needspace{8\baselineskip}

\begin{problem}[points=12]
  \noindent
  \begin{minipage}[t]{\dimexpr\linewidth-52mm-6mm\relax}
    \vspace{0pt}
    如图，在四边形 $ABCD$ 中，$\angle BAD=90^\circ$，$CD\perp BD$，$BD^2=AD\cdot BC$，点 $H$ 在边 $BD$ 上，且 $\angle DBC=\angle C+\angle BAH$。
    （1）求证：$AH=AD$；（2）过点 $B$ 作射线 $BE$ 交边 $CD$ 于点 $F$，交 $AD$ 的延长线于点 $E$，若 $\angle DBE=\angle ABD$，求证：$BD\cdot DE=EF\cdot CD$。
  \end{minipage}\hfill
  \begin{diagramcoltikz}{52mm}{}
  \includegraphics[width=\linewidth]{\detokenize{assets/Q009/prompt-figure.png}}
  \end{diagramcoltikz}
  \par\medskip
  \answerarea[72mm]
\end{problem}

% ---------- 第 10 题 (杨浦23) ----------

\needspace{8\baselineskip}

\begin{problem}[points=12]
  \noindent
  \begin{minipage}[t]{\dimexpr\linewidth-52mm-6mm\relax}
    \vspace{0pt}
    已知：如图，在矩形 $ABCD$ 中，点 $E$ 在边 $AB$ 的延长线上，过 $A$ 作 $AF\perp CE$，垂足为点 $F$，$CF=EF$，$AF$ 与边 $BC$ 交于点 $G$，联结 $BF$。
    （1）求证：$FB^2=FA\cdot FG$；（2）连接 $BD$ 与 $AG$ 交于点 $H$，如果 $AB=BC$，求证：$\dfrac{AH}{HG}=\dfrac{HF}{FG}$。
  \end{minipage}\hfill
  \begin{diagramcoltikz}{52mm}{}
  \includegraphics[width=\linewidth]{\detokenize{assets/Q010/manual-prompt-1-a6fdaceab9e934f1.png}}
  \end{diagramcoltikz}
  \par\medskip
  \answerarea[72mm]
\end{problem}

\end{document}